% NepalDTM-G: a national 1 m terrain model of Nepal by lidar-trained generative downscaling of NepalDTM v2 (8 m).
% Compile from this directory: tectonic paper_g.tex. Figures: make_figures.py (run from the repo root).
\documentclass[11pt]{article}

\usepackage[T1]{fontenc}
\usepackage[margin=1in]{geometry}
\usepackage{booktabs}
\usepackage{amsmath}
\usepackage{graphicx}
\usepackage{xurl}
\usepackage{xcolor}
\usepackage[section]{placeins}  % figures and tables never float past the end of their section
\usepackage[round,authoryear]{natbib}
\usepackage[hidelinks]{hyperref}
\graphicspath{{figures/}}
\newcommand{\fig}[4][0.42]{\begin{figure}[!htbp]\centering\includegraphics[width=0.95\textwidth,height=#1\textheight,keepaspectratio]{#2.pdf}\caption{#3}\label{#4}\end{figure}}
\hypersetup{%
  pdftitle={NepalDTM-G: A 1 m Digital Elevation Model of Nepal from Diffusion-Based Super-Resolution Trained on Lidar},
  pdfkeywords={Nepal, DEM, DTM, digital elevation model, 1 m, high resolution, super-resolution, diffusion model, deep learning, lidar, Himalaya},
  pdfauthor={Achyut Paudel}%
}

\newcommand{\gtitle}{NepalDTM-G: A 1~m Digital Elevation Model of Nepal\\[2pt] from Diffusion-Based Super-Resolution\\[2pt] Trained on Lidar}
\newcommand{\orcidid}{0000-0002-1459-3631}
\newcommand{\affil}{Independent Researcher}
\newcommand{\email}{\href{mailto:achyut.pdl@gmail.com}{achyut.pdl@gmail.com}}
\newcommand{\journalnote}{This manuscript has not been submitted to a peer-reviewed journal.}
\title{\gtitle}
\author{Achyut Paudel\\
\small \affil\\
\small \email\quad ORCID: \url{https://orcid.org/\orcidid}}
\date{November 2026}

\setlength{\emergencystretch}{2em}
\begin{document}

\begin{titlepage}
\centering
\vspace*{2cm}
{\LARGE \gtitle\par}
\vspace{1.5cm}
{\large Achyut Paudel\par}
\vspace{4pt}
\affil\par
Email: \email\par
ORCID: \url{https://orcid.org/\orcidid}\par
\vspace{1.5cm}
\begin{minipage}{0.82\textwidth}
\centering
This manuscript is a \textbf{non-peer-reviewed preprint} submitted to EarthArXiv.\\[4pt]
\journalnote\\[12pt]
Code: \url{https://github.com/achyut/NepalDTM} (branch \texttt{g})\\
Viewer: \url{https://nepaldtm.achyutpdl.com/map?v=g}\\
Input model (NepalDTM v2): \url{https://doi.org/10.31223/X5B23B}; data \url{https://doi.org/10.5281/zenodo.24327134}\\
NepalDTM v1: \url{https://doi.org/10.31223/X5DR5M}\\[12pt]
Comments and feedback are welcome: \email
\end{minipage}
\vfill
{\small November 2026\par}
\end{titlepage}
\setcounter{page}{1}
\maketitle

\begin{abstract}
No open digital elevation model (DEM) of Nepal exists at 1~m. The finest national measured model is NepalDTM v2 (M), a bare-earth
model on an 8~m grid. NepalDTM-G (G) adds 1~m detail to M without changing it at 8~m. G is a deep-learning super-resolution model, a conditional denoising
diffusion model trained on 75{,}521 lidar patches from the Swiss Alps and three United States mountain ranges, each
degraded with an error model fitted to M, so that training inputs look like the real product rather than clean block
averages. A smooth mass-preserving projection, applied after the samples are combined and the windows blended, makes
every $8\times8$~m block of G average to its M cell within 2~mm. On held-out lidar patches the final model is 3.3\% more
accurate than bicubic interpolation and matches the true slope field better (correlation 0.70 against 0.63); ten
sampling steps with two samples give the same result at a tenth of the cost. Against held-out ICESat-2 photons in two
test areas of Chitwan, G is 1--2\% worse than M. We show that this is a limit of the test, not of G: with the few metres
of horizontal error of the photons, even the true 1~m surface scores worse than bicubic, so point photons can check G's
8~m content, which equals M, but not its detail. G's detail is generated, not measured, and it is published as a
product separate from M.
\end{abstract}

\noindent\textbf{Keywords:} Nepal; digital elevation model (DEM); digital terrain model (DTM); 1~m; high resolution; super-resolution; diffusion model; deep learning; lidar; Himalaya

\section{Introduction}
NepalDTM is an open national terrain model of Nepal on an 8~m grid~\citep{paudel2026nepaldtm}. Its second release (v2,
called M here) adjusts the NASA High Mountain Asia 8~m stereo strips~\citep{shean2017hma,shean2020hma} to ICESat-2 ATL03
ground photons~\citep{markus2017icesat2,neumann2019atl03} and scores 5.26~m RMSE against held-out photons
nationally~\citep{paudel2026nepaldtmv2}. Many uses want finer detail: drainage on terraced slopes, gully heads, river
banks, road cuts and web maps at street zooms. No open source measures Nepal at 1~m, so a 1~m product has to be
generated, with fine structure learned from airborne lidar elsewhere and placed on the measured 8~m surface.

Such a product can be mistaken for measurement, so G obeys two rules. Averaged back to 8~m, G equals M, so it cannot
change M's measured accuracy. And its 1~m detail is tested against lidar, the only data that measure terrain at that
scale; ICESat-2 photons, which validate M, are also scored, and Section~\ref{sec:limits} shows what they can and cannot
tell about detail below 8~m. Within those rules G's job is to make the relief below 8~m realistic.

Learned super-resolution of DEMs is established~\citep[e.g.][]{chen2016demsr}, and diffusion models give the most
realistic image super-resolution~\citep{saharia2023sr3}. The difficulty here is that the input is a real product with
real errors. In an earlier NepalDTM experiment, a residual network trained on clean 8~m block means of lidar beat bicubic
on clean held-out inputs (1.89 against 2.25~m RMSE) but gave no gain on the 8~m product itself (6.83 against 6.86~m): it
read the product's stereo noise as terrain. G is therefore trained on lidar degraded to look like M.
Figure~\ref{fig:location} shows the areas used.
\fig[0.36]{fig06_location_map}{NepalDTM v2 over Nepal, with the area where the error model of M was fitted (Langtang)
and the two test areas in Chitwan District.}{fig:location}

\section{Data}
\subsection{Error model of M}
Against the independent 2~m WorldView/GeoEye composite at Langtang~\citep{shean2026langtang}, the error of M is mildly
blurred (Gaussian $\sigma=0.45$~cells), spatially correlated (e-folding length 11~m) and grows with slope; voids, filled
from coarser data, rise from 0.2\% of cells below $10^\circ$ to 9.2\% above $50^\circ$ (Table~\ref{tab:errmodel}). A
degradation function applies these statistics to clean lidar: block-average to 8~m, blur, add correlated noise scaled
per slope bin, and replace a slope-dependent share of cells with smoothed, noisier values. It reproduces the NMAD of M in
every slope bin within 3\% (Figure~\ref{fig:errmodel}); this is partly by construction, and whether the inputs are
realistic is tested on Nepal itself (Section~\ref{sec:nepal}).
\begin{table}[!htbp]
\centering\small
\caption{Stereo error of M at Langtang by slope (NMAD, m), share of void cells, and the same NMAD on degraded lidar.}
\label{tab:errmodel}
\begin{tabular}{lrrrrrr}
\toprule
Slope & $0$--$10^\circ$ & $10$--$20^\circ$ & $20$--$30^\circ$ & $30$--$40^\circ$ & $40$--$50^\circ$ & $\geq50^\circ$ \\
\midrule
NMAD, M at Langtang & 0.39 & 0.50 & 0.65 & 0.88 & 1.29 & 2.11 \\
Void share & 0.2\% & 0.7\% & 1.7\% & 2.7\% & 4.9\% & 9.2\% \\
NMAD, degraded lidar & 0.39 & 0.50 & 0.66 & 0.89 & 1.31 & 2.16 \\
\bottomrule
\end{tabular}
\end{table}
\fig[0.3]{fig09_error_model}{Stereo error of M by slope at Langtang and on degraded lidar (bars, left axis), and the
share of void cells (line, right axis).}{fig:errmodel}

\subsection{Training lidar}
We cut 75{,}521 gap-free $256\times256$~m patches at 1~m (Table~\ref{tab:data}, Figure~\ref{fig:sites}): 39{,}150 from
swissALTI3D~\citep{swisstopo_alti3d} (0.5~m, averaged to 1~m) and 36{,}371 from USGS 3DEP 1~m lidar~\citep{usgs3dep} in
three settings chosen to span Nepal's mountain terrain. Every 20th patch is held out. None of the areas is in Nepal, and
none resembles the Terai floodplain or the Churia (Siwalik) hills closely.
\begin{table}[!htbp]
\centering\small
\caption{Training patches ($256\times256$ at 1~m).}
\label{tab:data}
\begin{tabular}{llr}
\toprule
Source & Setting & Patches \\
\midrule
swissALTI3D 0.5~m & Valais and Bernese Alps & 39{,}150 \\
USGS 3DEP 1~m & San Juan Mountains, Colorado (high, rocky) & 13{,}332 \\
USGS 3DEP 1~m & Cascades and Olympics, Washington (steep forest) & 11{,}643 \\
USGS 3DEP 1~m & Great Smoky Mountains (forested mid-hills) & 11{,}396 \\
\midrule
 & total & 75{,}521 \\
\bottomrule
\end{tabular}
\end{table}
\fig[0.25]{fig07_training_sites}{One training patch per site, at the 80th percentile of relief within the site.}{fig:sites}

\section{The AI model}
\label{sec:model}
G is a conditional denoising diffusion model~\citep{ho2020ddpm} (Figure~\ref{fig:pipeline}, Table~\ref{tab:hyper}). From
an 8~m window of $32\times32$ cells it predicts the 1~m residual $r = z_{1\,\mathrm{m}} - \mathrm{bic}(x_{8\,\mathrm{m}})$,
the detail that bicubic upsampling misses, so it never has to learn absolute height. The denoiser is a
U-Net~\citep{ronneberger2015unet} of 14~M parameters conditioned on the bicubic upsampling and its gradients, with a cosine
noise schedule~\citep{nichol2021improved} and $v$-prediction~\citep{salimans2022progressive}. Each training patch is
degraded afresh with the error model of M, rotated or flipped, and the loss minimised with
AdamW~\citep{loshchilov2019adamw}; an exponential moving average of the weights is used for inference. Training ran for
150{,}000 steps on one laptop, a MacBook Pro with an Apple M4 Max (32-core GPU, 36~GB unified memory) using PyTorch~2.15
with its Metal (MPS) backend, at about 1.3~s per step (55--80 hours); the same machine ran all inference. The loss falls
quickly over the first 2{,}000 steps and slowly thereafter (Figure~\ref{fig:loss}).
\fig[0.3]{fig10_pipeline}{G in training (top) and inference (bottom).}{fig:pipeline}
\begin{table}[!htbp]
\centering\small
\caption{Model, training and inference settings of G.}
\label{tab:hyper}
\begin{tabular}{ll}
\toprule
Input / output & $32\times32$ cells at 8~m / $256\times256$ cells at 1~m \\
Target & 1~m residual from bicubic upsampling, in units of 2~m \\
Conditioning & bicubic upsampling (in 50~m units) and its $x$, $y$ gradients \\
Denoiser & U-Net, 4 levels (64, 128, 192, 256 channels), 14.0~M parameters \\
Diffusion & 1{,}000 steps, cosine schedule, $v$-prediction \\
Optimiser & AdamW, learning rate $2\times10^{-4}$, 2{,}000-step warm-up, weight decay $10^{-4}$ \\
Batch, augmentation & 12 patches; rotations by $90^\circ$ and flips; degradation redrawn each time \\
Weights used & exponential moving average, decay 0.999 \\
Sampling & DDIM, 10--50 steps, median of 2--4 samples (Section~\ref{sec:cost}) \\
Projection & smooth block-mean correction, 8 passes, per sample and after blending \\
Windows & 256~m every 192~m (64~m overlap), raised-cosine blending \\
Hardware & one laptop GPU (Apple M4 Max, 36~GB), PyTorch MPS backend \\
\bottomrule
\end{tabular}
\end{table}
\fig[0.3]{fig01_training_loss}{Training loss, as logged (light) and smoothed (dark). Dotted lines mark the checkpoints
evaluated on held-out patches. The run was paused and resumed at 24{,}000 steps.}{fig:loss}

\section{Inference}
\label{sec:inference}
\subsection{Sampling and the ensemble median}
DDIM sampling~\citep{song2021ddim} turns a noise field into one plausible 1~m surface; samples differ only in detail the
8~m input does not fix. We keep the per-pixel median of $K$ samples, which retains what they agree on. One sample is no
more accurate than bicubic; the median of four is (Section~\ref{sec:heldout}).

\subsection{Mass-preserving projection}
\label{sec:projection}
Each $8\times8$~m block of G must average to its M cell. Shifting each block by its own constant does this exactly but
leaves a step at every cell edge, a faint grid with twice the normal roughness there (Figure~\ref{fig:seams}). Instead the
block-mean error is upsampled bicubically and added, eight times, each pass cutting it about threefold. This is done for
every sample and again for each blended tile, because the median and the window blending undo it: without the final pass
block means were off by 0.17--0.25~m RMSE (up to 2~m); with it they are within 1.5~mm of M, with no grid and under a
second of extra time. G averaged to 8~m therefore returns M and cannot change its accuracy.
\fig[0.27]{fig05_seam_diagnostic}{Roughness (mean absolute second difference along rows) by column position within the
8~m cell, relative to its median. The block-constant projection (Chitwan) peaks at the cell edges; the smooth projection
(Churia) is flat. Bicubic (dotted) has its own smooth periodicity from the interpolation kernel.}{fig:seams}

\subsection{From windows to a province}
Windows of 256~m sit every 192~m on one lattice of the national 8~m grid and are blended with a raised-cosine weight; voids
in M are filled from the nearest valid cell and masked in the output. A province is processed in resumable 10~km tiles,
each computing every window that touches it and cropped to the province. Each window's noise is seeded by its lattice
position and DDIM is deterministic, so a window shared by two tiles is identical in both and tiles join without seams.

\subsection{Cost and cheaper sampling}
\label{sec:cost}
Run time scales with network evaluations, steps $\times$ samples; faster GPUs shorten it but not the ratios between
settings. Table~\ref{tab:speed} and Figure~\ref{fig:speed} score cheaper settings of the final model on 256 held-out
patches against a tolerance fixed in advance (0.01~m RMSE, 0.005 slope correlation):
\begin{itemize}\setlength\itemsep{1pt}
\item 20 steps $\times$ 4 samples is the best setting (gain over bicubic 3.7\%, slope correlation 0.706) and 2.7 times
faster than 50 steps: fewer steps cost nothing.
\item 10 steps $\times$ 2 samples matches 50 steps (3.4\% against 3.3\%; slope correlation 0.697 against 0.700) at 11 times
less cost.
\item Half precision gave identical scores; its speed-up varied between runs (none at 110{,}000 steps, 17--28\% at
150{,}000), so it is not used.
\item Sampling noise moves scores by 0.01--0.02~m, which explains $20\times2$ scoring below $10\times2$; the small edge of
$20\times4$ over $10\times2$ is consistent in both measures.
\end{itemize}
Bagmati Province (20{,}300~km$^2$) is generated with $10\times2$. Its 261 tiles of 10~km hold 0.73~million windows
including the tile margins, and the run measures 0.53~s per window including reading, blending and writing, about
4.5 days on the laptop; $20\times4$ would take about 19 days and $50\times4$ about 50. With
LERC compression~\citep{esri_lerc} at 1~cm maximum error, as for M, the 1~m surface takes 1.1~MB per km$^2$: about 22~GB
for Bagmati and 160~GB for Nepal.
\begin{table}[!htbp]
\centering\small
\caption{Sampling settings of the final model (150{,}000 steps) on 256 held-out patches. Gain is the RMSE reduction
against bicubic (1.294~m). Speed-up is against 50 steps with four samples; days for Bagmati scale the measured 0.53~s per
window of $10\times2$.}
\label{tab:speed}
\begin{tabular}{lrrrrr}
\toprule
Steps $\times$ samples & RMSE (m) & Gain & Slope corr. & Speed-up & Bagmati (days)\\
\midrule
$50\times4$ (reference) & 1.251 & 3.3\% & 0.700 & 1.0$\times$ & 50\\
$20\times4$ & 1.246 & 3.7\% & 0.706 & 2.7$\times$ & 19\\
$20\times2$ & 1.261 & 2.6\% & 0.685 & 5.9$\times$ & 9\\
$10\times2$ (used) & 1.250 & 3.4\% & 0.697 & 11.2$\times$ & 4.5\\
$20\times1$ & 1.278 & 1.2\% & 0.675 & 11.8$\times$ & 4\\
\bottomrule
\end{tabular}
\end{table}
\fig[0.3]{fig11_sampling_speed}{Cost against quality of the sampling settings in Table~\ref{tab:speed}. Shading marks
the tolerance around 50 steps with four samples; the dashed line is bicubic.}{fig:speed}

\section{Results}
\subsection{Held-out lidar patches}
\label{sec:heldout}
Held-out patches are degraded with the error model, using a fixed seed per patch, and the 1~m outputs are compared with
the true lidar (Table~\ref{tab:heldout}, Figure~\ref{fig:heldout}); from 50{,}000 steps the sampling noise is seeded too.
A single G sample is no better than bicubic, while the median of four is 2.8--2.9\% better from 30{,}000 to 110{,}000
steps and 3.3\% better at 150{,}000, the first evaluation with the final projection after the median. The clearer gain is
in shape: the slope map of G correlates with the true slope map at 0.69--0.70, against 0.63 for bicubic, so the added
relief is partly right, not only plausible.
\begin{table}[!htbp]
\centering\small
\caption{Held-out lidar patches: RMSE (m) against the true 1~m surface and correlation of slope maps with the true slope
map. Rows 4{,}000 and 12{,}000 share 64 patches; the later rows share 256. From 150{,}000 steps the median is projected
onto M after it is taken (Section~\ref{sec:projection}).}
\label{tab:heldout}
\begin{tabular}{rrrrrrrr}
\toprule
 & & \multicolumn{4}{c}{RMSE (m)} & \multicolumn{2}{c}{Slope-map corr.} \\
\cmidrule(lr){3-6}\cmidrule(lr){7-8}
Step & Patches & 8~m nearest & Bicubic & G, 1 sample & G, median of 4 & Bicubic & G \\
\midrule
4{,}000 & 64 & 1.844 & 1.186 & 1.185 & 1.159 & 0.628 & 0.639 \\
12{,}000 & 64 & 1.844 & 1.186 & 1.198 & 1.158 & 0.628 & 0.632 \\
30{,}000 & 256 & 1.994 & 1.294 & 1.300 & 1.258 & 0.627 & 0.691 \\
50{,}000 & 256 & 1.994 & 1.294 & 1.297 & 1.259 & 0.627 & 0.693 \\
88{,}000 & 256 & 1.994 & 1.294 & 1.295 & 1.257 & 0.627 & 0.696 \\
110{,}000 & 256 & 1.994 & 1.294 & 1.297 & 1.258 & 0.627 & 0.696 \\
150{,}000 & 256 & 1.994 & 1.294 & 1.293 & 1.251 & 0.627 & 0.700 \\
\bottomrule
\end{tabular}
\end{table}
\fig[0.3]{fig02_heldout_patches}{Held-out lidar patches by checkpoint (steps, and number of patches). Left: RMSE against
the true 1~m surface. Right: RMSE gain of the G median over bicubic (bars, left axis) and slope-map correlation (points,
right axis).}{fig:heldout}

\subsection{Nepal: Chitwan floodplain and Churia hills}
\label{sec:nepal}
The final model ($10\times2$) was applied to two areas of Chitwan District unlike any training terrain: the Rapti
floodplain near Sauraha (84.40--84.50$^\circ$E, 27.55--27.65$^\circ$N; 155--228~m) and forested ridges of the Churia range
(84.45--84.50$^\circ$E, 27.46--27.51$^\circ$N; about 250--650~m). Both are scored against held-out ICESat-2 photons from
tracks never used to build M (Table~\ref{tab:nepal}).

G adds 9--22\% more sub-grid relief than bicubic. On the floodplain it draws channels, banks and terrace edges that
bicubic blurs (Figure~\ref{fig:chitwan}); in the hills it sharpens ridges and gully heads (Figure~\ref{fig:churia}), and
the added relief follows the drainage network rather than appearing as uniform texture (Figure~\ref{fig:churiamaps}).
Against the photons G is 2.0\% worse than M on the floodplain (2.02 against 1.98~m) and 2.0\% worse in the hills (11.87
against 11.63~m; 95\% interval of the difference 0.05--0.39~m from a block bootstrap along the tracks), while bicubic
matches M in both. The median error of G equals that of M, so the difference is added scatter, not bias;
Section~\ref{sec:limits} shows why point photons penalise any 1~m detail. The hills box also shows a limit of the
approach: M there is Copernicus fill under forest, about 7~m above the photons, and G cannot remove an error it is
given. A straight east--west step in the lower third of the box is in M itself.
\begin{table}[!htbp]
\centering\small
\caption{Chitwan test areas: error against held-out ICESat-2 ground photons (m), and the standard deviation of the 1~m
surface about its 8~m cell value (sub-grid relief, m).}
\label{tab:nepal}
\begin{tabular}{lrrrrr}
\toprule
Area & Product & RMSE & NMAD & Median & Sub-grid relief \\
\midrule
Floodplain (117{,}703 photons) & M, 8~m & 1.98 & 0.52 & $+0.13$ & -- \\
 & Bicubic, 1~m & 1.97 & 0.51 & $+0.13$ & 0.58 \\
 & G, 1~m & 2.02 & 0.53 & $+0.13$ & 0.71 \\
\midrule
Churia hills (1{,}610 photons) & M, 8~m & 11.63 & 7.93 & $+7.28$ & -- \\
 & Bicubic, 1~m & 11.65 & 8.44 & $+7.10$ & 1.19 \\
 & G, 1~m & 11.87 & 8.51 & $+7.25$ & 1.30 \\
\bottomrule
\end{tabular}
\end{table}
\fig{fig03_chitwan_floodplain}{Chitwan floodplain: hillshades of M (8~m, shown as blocks), bicubic and G at 1~m, over
1.5~km (top) and 400~m (bottom, red square); final model, $10\times2$.}{fig:chitwan}
\fig{fig04_churia_hills}{Churia hills, as Figure~\ref{fig:chitwan}; 2~km (top) and 400~m (bottom) views.}{fig:churia}
\fig[0.36]{fig08_churia_detail_and_photons}{Churia hills. Left: G minus bicubic, the relief G adds below 8~m. Right: G
minus held-out ICESat-2 ground photons along two tracks, over shaded relief of G.}{fig:churiamaps}

\subsection{What point photons can and cannot test}
\label{sec:limits}
Two tests explain the photon result. First, on 400 held-out lidar patches we placed 200 points per patch, took the true
height there, and read two surfaces at the point displaced by a random horizontal error: the true 1~m surface itself
and bicubic of its 8~m means (Table~\ref{tab:limits}). With no displacement the truth is perfect; with 2~m it only ties
with bicubic, and from 3~m on it scores worse. A sharp surface, even the right one, is punished for small position
errors that a smooth one absorbs, and ICESat-2 photon positions carry errors of a few metres. Second, in both Nepal
areas we scaled G's added detail, bicubic $+\,a\,(\mathrm{G}-\mathrm{bicubic})$: photon RMSE rises steadily with $a$, and
even a quarter of the detail scores worse than none. Photons therefore check G's 8~m content, which equals M by
construction, but cannot tell right detail from wrong; lidar can (Section~\ref{sec:heldout}), and lidar of Nepal would
settle it (Section~\ref{sec:future}).
\begin{table}[!htbp]
\centering\small
\caption{Point tests of 1~m detail (RMSE, m). Left: held-out lidar, true 1~m surface and bicubic of its 8~m means, read at
points displaced by Gaussian horizontal error $\sigma$. Right: held-out photons in Chitwan, bicubic $+\,a\,$(G $-$
bicubic).}
\label{tab:limits}
\begin{tabular}{rrr@{\qquad}rrr}
\toprule
$\sigma$ (m) & True 1~m & Bicubic & $a$ & Floodplain & Churia hills \\
\midrule
0 & 0.000 & 0.407 & 0 (bicubic) & 1.974 & 11.65 \\
1 & 0.582 & 0.677 & 0.25 & 1.983 & 11.70 \\
2 & 1.176 & 1.187 & 0.5 & 1.994 & 11.75 \\
3 & 1.722 & 1.707 & 0.75 & 2.007 & 11.81 \\
4 & 2.295 & 2.261 & 1 (G) & 2.023 & 11.87 \\
6 & 3.336 & 3.295 & M, 8~m & 1.983 & 11.63 \\
\bottomrule
\end{tabular}
\end{table}

\subsection{Comparison with other terrain models}
\label{sec:compare}
Table~\ref{tab:compare} scores G with the public global models and both NepalDTM releases on the same held-out photons
(only photons where every product has a value). Copernicus GLO-30~\citep{esa2021copernicus} and
FABDEM~\citep{hawker2022fabdem} are resampled bilinearly to the 8~m grid; GEM-15~\citep{gem15} is its 15~m DTM, converted
from EGM96~\citep{lemoine1998egm96} to EGM2008~\citep{pavlis2012egm2008} like the photons, and
GEDTM30~\citep{ho2025gedtm30} is resampled like Copernicus. On the floodplain M and bicubic are 8--9\% better than v1 and
59\% better than Copernicus, FABDEM is 0.10~m better still, and G is 2\% behind M for the reason of
Section~\ref{sec:limits}. In the Churia hills v1 is best by about 1~m, because v2 there is Copernicus fill with a canopy
residual of about $+7$~m. GEM-15 sits 16--31~m above the photons in both boxes.
\begin{table}[!htbp]
\centering
\caption{Terrain models on the same held-out ICESat-2 ground photons in the Chitwan test boxes (m; error is model minus
photon). Best RMSE per area in bold. G is the final model, $10\times2$.}
\label{tab:compare}
\small\setlength{\tabcolsep}{4pt}
\begin{tabular}{lrrrrrrr}
\toprule
 & & \multicolumn{3}{c}{Floodplain (101{,}604 photons)} & \multicolumn{3}{c}{Churia hills (1{,}610 photons)} \\
\cmidrule(lr){3-5}\cmidrule(lr){6-8}
Model & Grid & RMSE & NMAD & Median & RMSE & NMAD & Median \\
\midrule
Copernicus GLO-30 & 30~m & 4.91 & 0.63 & $+0.48$ & 22.37 & 7.53 & $+21.96$ \\
FABDEM V1-2 & 30~m & \textbf{1.92} & 0.55 & $+0.47$ & 11.89 & 8.13 & $+8.05$ \\
GEDTM30 & 30~m & 2.00 & 0.62 & $+0.02$ & 12.44 & 7.83 & $+9.09$ \\
GEM-15 DTM & 15~m & 17.13 & 1.54 & $+15.99$ & 32.05 & 6.90 & $+30.85$ \\
NepalDTM v1~\citep{paudel2026nepaldtm} & 8~m & 2.20 & 0.45 & $+0.13$ & \textbf{10.64} & 8.99 & $+3.44$ \\
NepalDTM v2 (M)~\citep{paudel2026nepaldtmv2} & 8~m & 2.02 & 0.44 & $+0.11$ & 11.63 & 7.93 & $+7.28$ \\
Bicubic of M & 1~m & 2.01 & 0.44 & $+0.11$ & 11.65 & 8.44 & $+7.10$ \\
G & 1~m & 2.06 & 0.46 & $+0.11$ & 11.87 & 8.51 & $+7.25$ \\
\bottomrule
\end{tabular}
\end{table}

\section{Discussion and limitations}
\paragraph{What G is.} G is a prior on how terrain looks at 1~m, conditioned on M. Where many 1~m surfaces fit the
8~m data equally well, it picks a realistic one, and the median of samples keeps the features they agree on. The
slope-map gain on held-out lidar shows that part of this structure is real, but nothing in Nepal can yet say which 1~m
features are: ICESat-2 photons cannot judge 1~m detail (Section~\ref{sec:limits}), and there is no open 1~m lidar of
Nepal. G must not be used where measured 1~m heights are required.

\paragraph{Point accuracy.} Read at a point, G is about 2\% less accurate than M or bicubic in the Chitwan tests, the
price of any sharp surface under position error. For uses that need heights at points (levelling, volumes over small
areas) M or its bicubic upsampling are the better choice; G is for the shape of the terrain: slopes, drainage, breaklines
and visualisation.

\paragraph{Inherited error.} G inherits all of M's error at 8~m, including canopy residuals under forest and steps
between source areas; it adds detail to the surface it is given.

\paragraph{Training domain.} The training lidar is alpine and mid-latitude mountain terrain; Nepal's floodplains,
terraces and the soft rock of the Churia have no counterpart in it (Section~\ref{sec:future}).

\paragraph{Release.} G is first released for Bagmati Province with the setting chosen in Section~\ref{sec:cost} (in
production at the time of writing), shown on
the web map at zooms 15--17 as an option separate from the measured models. A national release depends on its cost and
on a comparison against the 2~m composites at Langtang~\citep{shean2026langtang} and
Chamoli~\citep{bhushan2021chamoli}.

\section{Future work: national lidar from the Survey Department}
\label{sec:future}
The Survey Department of Nepal has flown airborne lidar over about 20{,}000~km$^2$ of the western Terai and Chure, from
Chitwan to Kanchanpur, for a 1~m terrain model of 0.25~m accuracy, 0.5~m contours and orthophotos~\citep{joshi2020lidar}.
To our knowledge the data are not openly released; using them would need an agreement with the Department. They would:
\begin{itemize}\setlength\itemsep{1pt}
\item \textbf{Test G at 1~m in Nepal}, on held-out lidar blocks, turning the go/no-go for a national G into a measurement
(the Chitwan test areas may lie inside the survey).
\item \textbf{Train G on Nepali terrain}: floodplain, terraces and Churia hills, which the training set lacks; the error
model of M, now fitted only at Langtang, would be refitted there too.
\item \textbf{Correct M}, by measuring its error cell by cell, including the forest canopy residual (about $+7$~m in the
Churia box), where ICESat-2 ground photons are sparse.
\item \textbf{Replace G where flown}: measured 1~m heights inside the survey, M and G outside, blended at the boundary.
\item \textbf{Validate all NepalDTM releases} at a scale ICESat-2 cannot.
\end{itemize}

\section{Data and code availability}
\begin{sloppypar}
Code: \url{https://github.com/achyut/NepalDTM}, branch \texttt{g}, folder \texttt{code/g/}: error model
(\texttt{error\_model.py}), training patches (\texttt{make\_patches.py}), model and training (\texttt{g\_model.py}),
inference (\texttt{g\_apply.py}, \texttt{g\_region.py}, \texttt{g\_tiles.py}), sampling benchmark (\texttt{g\_bench.py}),
model comparison (\texttt{g\_compare.py}) and photon tests (\texttt{g\_photon\_limits.py}).
\end{sloppypar} Input model M: NepalDTM v2~\citep{paudel2026nepaldtmv2}, data at
\url{https://doi.org/10.5281/zenodo.24327134}~\citep{paudel2026nepaldtmv2data}. G test tiles:
\url{https://nepaldtm.achyutpdl.com/map?v=g}. The model weights and the Bagmati release will be deposited on Zenodo.

\FloatBarrier
\bibliographystyle{plainnat}
\bibliography{references}

\end{document}
